%% ===========================================================
%% Copernicus Publications LaTeX manuscript
%% Requires copernicus.cls / copernicus.bst, available from
%% https://publications.copernicus.org/for_authors/manuscript_preparation.html
%% Replace "gmd" below with the abbreviation of the target journal.
%% ===========================================================
\documentclass[gmd, manuscript]{copernicus}

\begin{document}

\title{Automatic Camera Animation for Met.3D}

%% --- authors -------------------------------------------------
\Author[1]{}{Arpana}
\Author[1]{}{Samarth}
\Author[1]{Suraj}{Kumar}

\affil[1]{Department of Applied Computer Science, Fulda University of
Applied Sciences, Fulda, Germany}

\correspondence{xyz (email@example.com)}

\runningtitle{Automatic camera animation for Met.3D}
\runningauthor{S. Kumar et al.}

\received{}
\pubdiscuss{}
\revised{}
\accepted{}
\published{}

\firstpage{1}

\maketitle

%% ===========================================================
\begin{abstract}
Three-dimensional visualisation of weather forecast data requires a
camera path that shows the phenomena of interest clearly. Such paths are
normally created by hand for each case. We present a method that
generates camera animations automatically from detected and tracked
atmospheric features, using atmospheric rivers as the target phenomenon.
Given binary feature masks and per-timestep tracking records, the method
extracts a guide curve for each feature. From this curve it builds a set
of camera behaviours following the feature along its length, circling
it, and visiting several features in turn and writes them out as
camera sequences that run in the Met.3D visualisation framework.

Two design decisions proved important. First, points along the guide
curve must be ordered by walking the skeleton graph rather than by
sorting them along a coordinate axis, which otherwise causes the path to
double back on itself. Second, camera movement and forecast time must not
change at once: the camera stays still while the forecast advances, and
the forecast is held while the camera moves. We compare generated
sequences against hand-authored ones and describe the camera conventions
recovered from them.
\end{abstract}

%% ===========================================================
\introduction
%% ===========================================================

Numerical weather prediction produces large, four-dimensional datasets
that are difficult to inspect in their entirety. Three-dimensional
visualisation offers one route to understanding them, showing structures
such as fronts, cyclones and moisture transport in their spatial context
rather than as a stack of two-dimensional maps. Tools such as Met.3D make
this possible for interactive scientific work, and the same renderings
are increasingly used for communication in teaching material, in
public outreach, and in television weather broadcasts.

Any three-dimensional animation, however, requires a camera path: a
description of where the viewpoint sits at each moment, where it looks,
and how it moves. The choice of path is not incidental. A path that stays
too far away reveals little; one that moves too close loses geographic
context; one that moves erratically is tiring to watch. In practice these
paths are authored by hand, keyframe by keyframe, for each individual
case. This is acceptable for a single figure or a one-off animation, but
it does not scale to routine use, where forecasts are produced several
times a day and the features worth showing differ from one run to the
next.

This suggests an alternative: if the features of interest have already
been detected in the data, the camera path could be derived from them.
Detection and tracking of atmospheric features is by now well
established, and produces exactly the information a camera would need
where a feature is, what shape it has, how large it is, and how it moves
over time. The question is how to turn that information into camera
motion that is both informative and pleasant to watch.

We address this question for atmospheric rivers: long, narrow filaments
of concentrated water vapour that transport a substantial fraction of the
moisture reaching the mid-latitudes, and which are of direct forecasting
interest because of their association with heavy precipitation at
landfall. Their elongated geometry makes them a natural test case, since
it admits an obvious camera behaviour following the feature along its
length while their tendency to split, merge and deform over their
lifetime makes that behaviour non-trivial to realise.

The contribution of this work is a method for generating camera
animations automatically from detected and tracked atmospheric-river
features, together with an account of the design decisions that determine
whether the resulting animations are usable. We describe how a guide
curve is extracted from each feature, how several distinct camera
behaviours are parameterised from that curve, and how camera motion is
coordinated with the advance of forecast time. We also report the camera
conventions of the target visualisation framework, which we recovered
empirically from hand-authored sequences, since these proved to be the
difference between paths that keep the feature in view and paths that do
not.

%% ===========================================================
%% Further sections to be added:
%%   \section{System overview}
%%   \section{Methodology}
%%   \section{Results}
%%   \section{Discussion}
%%   \conclusions
%% ===========================================================

%% ===========================================================
%% Back matter (required by Copernicus)
%% ===========================================================

\codeavailability{} %% TODO

\dataavailability{} %% TODO

\authorcontribution{} %% TODO: e.g. feature tracking and ranking were
%% developed by ...; skeleton extraction and camera animation by ...

\competinginterests{The authors declare that they have no conflict of
interest.}

\begin{acknowledgements}
%% TODO
\end{acknowledgements}

\bibliographystyle{copernicus}
\bibliography{references.bib}

\end{document}
